% FORMAT AND PACKAGES
% {
\documentclass[a4paper, 11pt]{article}
\usepackage{a4wide,amssymb,epsfig,latexsym,multicol,array,hhline,fancyhdr}
\usepackage{vntex}
\usepackage{amsmath}
\usepackage{lastpage}
\usepackage[lined,boxed,commentsnumbered]{algorithm2e}
\usepackage{enumerate}
\usepackage{xcolor}
\usepackage{graphicx}							% Standard graphics package
\usepackage{array}
\usepackage{tabularx, caption}
\usepackage{multirow}
\usepackage{multicol}
\usepackage{rotating}
\usepackage{graphics}
\usepackage{geometry}
\usepackage{setspace}
\usepackage{epsfig}
\usepackage{tikz}
\usepackage{xfrac}
\usepackage{bm}
\usepackage{biblatex}
\usepackage[colorlinks]{hyperref}
% \usepackage[acronym,toc]{glossaries}
% \usepackage[symbols,nogroupskip,nonumberlist]{glossaries-extra}
\usepackage[
 sort=none,% no sorting or indexing required
 abbreviations,% create list of abbreviations
 symbols,% create list of symbols
 stylemods,style=list, % set the default glossary style
 nogroupskip, nonumberlist, nomain
]{glossaries-extra}


% FORMATTING
% {
\DeclareMathOperator{\arccot}{arccot}
\captionsetup[table]{name=Bảng}
\captionsetup[figure]{name=Hình}
\newenvironment{Description}{\list{}{%
    \let\makelabel\descriptionlabel    % this comes from the original description environment
    \setlength{\rightmargin}{\leftmargin}% this comes from the original quote environment
    \setlength{\labelwidth}{0pt}%          this is new
    }}{\endlist}

\addbibresource{citations.bib}
    
\hypersetup{urlcolor=blue,linkcolor=black,citecolor=black,colorlinks=true} 
\usetikzlibrary{arrows,snakes,backgrounds}
\definecolor{mathblue}{RGB}{0,114,188}
% \makeatletter  \def\m@th{\mathsurround\z@\color{mathblue}} \makeatother
% \everymath{\color{mathblue}}
% \setmathfont[Color=000000]{Arial}
%\usepackage{pstcol} 								% PSTricks with the standard color package
\newtheorem{theorem}{{\bf Theorem}}
\newtheorem{property}{{\bf Property}}
\newtheorem{proposition}{{\bf Proposition}}
\newtheorem{corollary}[proposition]{{\bf Corollary}}
\newtheorem{lemma}[proposition]{{\bf Lemma}}

\AtBeginDocument{\renewcommand{\listfigurename}{Danh sách hình ảnh}}
\AtBeginDocument{\renewcommand{\listtablename}{Danh sách bảng}}
\AtBeginDocument{\renewcommand*\contentsname{Mục lục}}
\AtBeginDocument{\renewcommand*\refname{References}}
%\usepackage{fancyhdr}

\setlength{\headheight}{40pt}
\pagestyle{fancy}
\fancyhead{} % clear all header fields
\fancyhead[L]{
 \begin{tabular}{rl}
    \begin{picture}(25,15)(0,0)
    \put(0,-8){\includegraphics[width=8mm, height=8mm]{images/hcmut.png}}
    %\put(0,-8){\epsfig{width=10mm,figure=hcmut.eps}}
   \end{picture}&
	%\includegraphics[width=8mm, height=8mm]{images/hcmut.png} & %
	\begin{tabular}{l}
		\textbf{\bf \ttfamily Trường Đại học Bách khoa - ĐHQG-HCM}\\
		\textbf{\bf \ttfamily Khoa Khoa học và Kỹ thuật Máy tính}
	\end{tabular} 	
 \end{tabular}
}
\fancyhead[R]{
	\begin{tabular}{l}
		\tiny \bf \\
		\tiny \bf 
	\end{tabular}  }
\fancyfoot{} % clear all footer fields
\fancyfoot[L]{\scriptsize \ttfamily Bài tập lớn Tính toán song song - \the\year}
\fancyfoot[R]{\scriptsize \ttfamily Trang {\thepage}/\pageref{LastPage}}
\renewcommand{\headrulewidth}{0.3pt}
\renewcommand{\footrulewidth}{0.3pt}

\setcounter{secnumdepth}{4}
\setcounter{tocdepth}{4}

\makeatletter
\newcounter {subsubsubsection}[subsubsection]
\renewcommand\thesubsubsubsection{\thesubsubsection .\@alph\c@subsubsubsection}
\newcommand\subsubsubsection{\@startsection{subsubsubsection}{4}{\z@}%
                                     {-3.25ex\@plus -1ex \@minus -.2ex}%
                                     {1.5ex \@plus .2ex}%
                                     {\normalfont\normalsize\bfseries}}
\newcommand*\l@subsubsubsection{\@dottedtocline{3}{10.0em}{4.1em}}
\newcommand*{\subsubsubsectionmark}[1]{}
% \def\m@th{\mathsurround\z@\color{mathblue}}
\makeatother
% }
% }

% ACRONYMS & SYMBOLS
% {
% \makeglossaries
\setabbreviationstyle{long-short}
\newabbreviation{ode}{ODE}{(First-Order) Ordinary Differential Equation}
\newabbreviation{ivp}{IVP}{Initial-Value Problem}
\newabbreviation{lte}{LTE}{Local Truncation Error}
\newabbreviation{ds}{DS}{Dynamical System}
\newabbreviation{fig}{Fig.}{Figure}
\newabbreviation{tab}{Tab.}{Table}
\newabbreviation{sys}{Sys.}{System of Equations}
\newabbreviation{eq}{Eq.}{Equation}
\newabbreviation{eg}{e.g.}{For Example}
\newabbreviation{ie}{i.e.}{That Is}
% \glsnoexpandfields
\glsxtrnewsymbol[description = {Set of natural numbers}]{natural}{\ensuremath{\mathbb{N}}}
\glsxtrnewsymbol[description = {Set of real numbers}]{real}{\ensuremath{\mathbb{R}}}
\glsxtrnewsymbol[description = {Set of positive real numbers}]{real_positive}{\ensuremath{\mathbb{R}^+}}

% }

% Code listing
\usepackage{listings}
% \usepackage{xcolor}

\definecolor{codegreen}{rgb}{0,0.6,0}
\definecolor{codegray}{rgb}{0.5,0.5,0.5}
\definecolor{codepurple}{rgb}{0.58,0,0.82}
\definecolor{backcolour}{rgb}{0.95,0.95,0.92}

\lstdefinestyle{mystyle}{
    backgroundcolor=\color{backcolour},   
    commentstyle=\color{codegreen},
    keywordstyle=\color{magenta},
    numberstyle=\tiny\color{codegray},
    stringstyle=\color{codepurple},
    basicstyle=\ttfamily\footnotesize,
    breakatwhitespace=false,         
    breaklines=true,                 
    captionpos=b,                    
    keepspaces=true,                 
    numbers=left,                    
    numbersep=5pt,                  
    showspaces=false,                
    showstringspaces=false,
    showtabs=false,                  
    tabsize=2
}

\lstset{style=mystyle}
% 

% Float [H]
\usepackage{float}
%

% DOCUMENT
\begin{document}

% TITLE PAGE
\begin{titlepage}
    \begin{center}
        ĐẠI HỌC QUỐC GIA THÀNH PHỐ HỒ CHÍ MINH \\
        TRƯỜNG ĐẠI HỌC BÁCH KHOA \\
        KHOA KHOA HỌC VÀ KỸ THUẬT MÁY TÍNH
    \end{center}

    \vspace{1cm}

    \begin{figure}[h!]
        \begin{center}
            \includegraphics[width=3cm]{images/hcmut.png}
        \end{center}
    \end{figure}

    \vspace{1cm}


    \begin{center}
        \begin{tabular}{c}
            % \multicolumn{1}{l}{\textbf{{\Large TÍNH TOÁN SONG SONG }}} \\
            ~~                                                                  \\
            \hline
            \\
            \multicolumn{1}{l}{\textbf{{\Large Tính toán song  song (CO5133)}}} \\
            \\
            \textbf{\textit{{\Huge Assignment 2}}}                              \\
            \\
            \hline
        \end{tabular}
    \end{center}

    \vspace{2cm}

    \begin{table}[h]
        \centering
        \begin{tabular}{rl}
            \hspace{3 cm}\textbf{Giảng viên hướng dẫn}:
                                          & TS. Diệp Thanh Đăng          \\
                                          &                              \\[10pt]
            \textbf{Sinh viên thực hiện}: & Lê Nguyễn Kim Khôi - 2311671 \\
                                          & Lê Trọng Tín - 2313452       \\
                                          & Phạm Trần Minh Trí - 2313622 \\
        \end{tabular}
    \end{table}

    \begin{center}
        {\footnotesize THÀNH PHỐ HỒ CHÍ MINH, \the\month/\the\year}
    \end{center}
\end{titlepage}

\pagebreak
\tableofcontents
\pagebreak

% Glossaries
% {}
% \printunsrtglossary[type={symbols}, title={Danh sách ký hiệu}]
% \printunsrtglossary[type={abbreviations}, title={Danh sách từ viết tắt}]
\pagebreak
\listoffigures
\listoftables
\pagebreak
\addcontentsline{toc}{section}{\listfigurename}
\addcontentsline{toc}{section}{\listtablename}

% }

% Member list
\section*{Danh sách thành viên \& phân công}
\addcontentsline{toc}{section}{Danh sách thành viên \& phân công}

\begin{center}
    \begin{table}[h]
        \centering
        \begin{tabular}{|>{\centering\arraybackslash}p{1cm}|>{\raggedright\arraybackslash}p{3.75cm}|>{\centering\arraybackslash}p{1.75cm}|>{\raggedright\arraybackslash}p{4.5cm}|>{\centering\arraybackslash}p{1.5cm}|}
            \hline
            \textbf{STT} & \textbf{Họ và tên} & \textbf{MSSV} & \textbf{Phân công} & \textbf{\% hoàn thành} \\
            \hline
            %%%%% Student 1 %%%%%%%%%%
            1            & Lê Nguyễn Kim Khôi & 2311671       & - Exercise 1       & 100\%                  \\
            \hline
            %%%%% Student 2 %%%%%%%%%%%
            2            & Lê Trọng Tín       & 2313452       & - Exercise 2       & 100\%                  \\
            \hline
            %%%%% Student 3 %%%%%%%%%%%
            3            & Phạm Trần Minh Trí & 2313622       & - Exercise 3       & 100\%                  \\
            \hline
        \end{tabular}
        \caption{\label{table1}Danh sách thành viên \& phân công}
    \end{table}
\end{center}


\pagebreak


\section{Introduction}

Hmm \cite{texbook}

\begin{figure}[H]
    \centering
    \includegraphics[width=0.5\linewidth]{images/hcmut.png}
    \caption{Red-black tree}
    \label{fig:enter-label}
\end{figure}

\begin{lstlisting}[language=Python, caption=Python example]
import numpy as np
    
def incmatrix(genl1,genl2):
    m = len(genl1)
    n = len(genl2)
    M = None #to become the incidence matrix
    VT = np.zeros((n*m,1), int)  #dummy variable
    
    #compute the bitwise xor matrix
    M1 = bitxormatrix(genl1)
    M2 = np.triu(bitxormatrix(genl2),1) 

    for i in range(m-1):
        for j in range(i+1, m):
            [r,c] = np.where(M2 == M1[i,j])
            for k in range(len(r)):
                VT[(i)*n + r[k]] = 1;
                VT[(i)*n + c[k]] = 1;
                VT[(j)*n + r[k]] = 1;
                VT[(j)*n + c[k]] = 1;
                
                if M is None:
                    M = np.copy(VT)
                else:
                    M = np.concatenate((M, VT), 1)
                
                VT = np.zeros((n*m,1), int)
    
    return M
\end{lstlisting}

\section{Các nền tảng share-memory được sử dụng}
\label{shared-mem}

\begin{enumerate}
    \item Nền tảng 1 (Laptop của Khôi):
          \begin{itemize}
              \item \textbf{CPU:}  Apple M2 (8-core CPU)
              \item \textbf{RAM:} 16 GB
              \item \textbf{OS:} Tahoe 26.6.2
          \end{itemize}
    \item Nền tảng 2 (Laptop của Tín):
          \begin{itemize}
              \item \textbf{CPU:} Intel Core i7 13620H, 10 core, 16 thread
              \item \textbf{RAM:} 16 GB
              \item \textbf{OS:} Windows 11 Home Single Language
          \end{itemize}
    \item Nền tảng 3 (Laptop của Trí):
          \begin{itemize}
              \item \textbf{CPU:} Intel Core i7 13650HX, 14 core, 20 thread
              \item \textbf{RAM:} 24 GB DDR5
              \item \textbf{OS:} Windows 11 - WSL Ubuntu
          \end{itemize}
\end{enumerate}
\section{Bài tập 1: The A* Algorithm in the Sequential World}
Khôi làm nhé


\section{Bài tập 2: The Concurrent Set Data Structure}
Tín làm nhé

\section{Bài tập 3: The A* Algorithm in the Parallel World}
Trí làm nhé

\section{Kết luận}



% for testing

% REFERENCES
\pagebreak
\nocite{*}
\printbibliography[
    heading=bibintoc,
    title={Tài liệu tham khảo}
]
\end{document}

